% META: {"id": "TEXP-MAT-CO-012", "chapitre": "TEXP-MATRICES-MARKOV", "type_objet": "corrige", "exercice_ref": "TEXP-MAT-EX-012", "capacites_codes": ["C5"], "sources_inspiration": [], "status": "generated"}
% BEGIN-VERIFY
% from sympy import Matrix, Rational
% T = Matrix([[Rational(7, 10), Rational(3, 10)], [Rational(4, 10), Rational(6, 10)]])
% pi0 = Matrix([[1, 0]])
% pi1, pi2, pi3 = pi0 * T, pi0 * T ** 2, pi0 * T ** 3
% assert pi1 == Matrix([[Rational(7, 10), Rational(3, 10)]])
% assert pi2 == Matrix([[Rational(61, 100), Rational(39, 100)]])
% assert pi3 == Matrix([[Rational(583, 1000), Rational(417, 1000)]])
% T2 = T ** 2
% assert T2 == Matrix([[Rational(61, 100), Rational(39, 100)],
%                      [Rational(52, 100), Rational(48, 100)]])
% assert T2[0, 1] == Rational(39, 100)
% assert pi2 == T2.row(0)
% END-VERIFY
\begin{corrige}{TEXP-MAT-EX-012}
\textbf{1.} $\pi_1=\pi_0T=\begin{pmatrix}0{,}7&0{,}3\end{pmatrix}$,
$\pi_2=\begin{pmatrix}0{,}61&0{,}39\end{pmatrix}$,
$\pi_3=\begin{pmatrix}0{,}583&0{,}417\end{pmatrix}$.

\textbf{2.} $T^2=\begin{pmatrix}0{,}61&0{,}39\\0{,}52&0{,}48\end{pmatrix}$.
Le coefficient $(1\,;\,2)$ vaut $0{,}39$ : c'est la probabilité d'être au
centre au bout de deux jours en partant de la gare.

\textbf{3.} $\pi_2$ est égale à la première ligne de $T^2$. C'est normal :
$\pi_0$ sélectionne la première ligne, et $\pi_2=\pi_0T^2$.
\end{corrige}
